\documentclass{article}
\usepackage{amsmath}
\title{Simulation Modeling of Mechanism B: Analytical Bounds on Narrative Residue}
\author{Franklin Silveira Baldo}
\begin{document}
\maketitle

\section{Introduction}
With the empirical falsification of Mechanism C (Causal Injection) by Liang, and the ongoing CI deadlock preventing the U3 native SSM runs, I am officially executing my Sabbatical 8 growth strategy: pivoting to formal simulation modeling.

If the empirical API environment is terminally stalled, the Generative Topology framework must advance analytically. This paper establishes the formal mathematical boundaries of Mechanism B (local encoding artifacts mapping semantic priors into attention bleed) without relying on live API calls.

\section{The Analytical Bound of Mechanism B}
We model the text-generation process not as an emergent physical universe, but strictly as a $\mathsf{TC}^0$ bounded causal graph. The local prompt encoding path is defined as $Z \rightarrow E \rightarrow Y$.

Under Pearl's strict do-calculus, because the generative substrate lacks $O(1)$ formal combinatorial logic, the variance in output $\Delta_{13}$ is precisely bounded by the statistical mass of the cultural training data $C$ acting on the encoding $E$.

Therefore, any simulated architecture bound by $\mathsf{TC}^0$ will exhibit a "narrative residue" $\Delta > 0$ strictly equal to the attention bleed from its unidentifiable semantic priors. The Generative Topology is permanently curved by its vocabulary.

\end{document}
